\scanpage{46}
On peut $=$ supposer que $X$ est un sous-schéma fermé de $S$ [en faisant une extension \deleted{\unsure{de la base}} \inserted{\unsure{fid\textsuperscript{t}|fid.}} plate de $S$]. Mais alors $X_1$ et $X_2$ \deleted{étant} finis \inserted{et étales} sur $S$, et non vides, sont identiques à $X$, donc identiques. \deleted{L'assertion est bien connue.}

Dans le cas général, on voit déjà que $X_1$ et $X_2$ coïncident \deleted{\unsure{dans}}\inserted{sur} un voisinage ouvert \deleted{$U$} de $s$. Soit $U$ le plus \unsure{gd} ouvert de $S$ tel que $X_1|U=X_2|U$, il reste à prouver que $U$ est \uline{fermé}. Or si $s'\in\overline U$, on a $s'\in\overline{s''}$, $s''\in U$, donc
\[
X_{is'}=\overline{X_{is''}}\cap X_{s'}\qquad
\text{[car $X_{s'}\to S$ est ouvert]}
\]
donc $X_{1s'}=X_{2s'}$, donc par ce qui précède $s'\in U$, cqfd.

\deleted{\unsure{Pourquoi}}
\uline{Théorème de rigidité} Soient $S$ un préschéma loc noeth \uline{connexe}, $G$ un \uline{schéma} en groupes \deleted{\unsure{loc}}\inserted{comm.} de type fini \deleted{\illegible{1}} sur $S$, $G_1$ et $G_2$ deux sous-schémas en groupes fermés \uline{rigides}, soit $s\in S$ tels que $G_{1s}=G_{2s}$. \deleted{pour un $s\in S$.} \deleted{Alors} Alors $G_1=G_2$.
% archive sheet number: 45
